\documentclass[11pt,a4paper]{article}

\usepackage[utf8]{inputenc}
\usepackage[T1]{fontenc}
\usepackage[french]{babel}
\usepackage{amsmath,amssymb}
\usepackage{geometry}
\usepackage{enumitem}
\usepackage{multicol}
\usepackage{tabularx}
\usepackage{booktabs}
\usepackage{xcolor}
\usepackage{mdframed}
\usepackage{lmodern}
\usepackage{microtype}
\usepackage{fancyhdr}
\usepackage{lastpage}
\usepackage{parskip}

\geometry{top=2cm, bottom=2.5cm, left=2.2cm, right=2.2cm}

% En-tête / pied de page
\pagestyle{fancy}
\fancyhf{}
\fancyhead[L]{\small\textit{Mathématiques - Analyse}}
\fancyhead[R]{\small\textit{Contrôle continu}}
\fancyfoot[C]{\small Page \thepage\ / \pageref{LastPage}}
\renewcommand{\headrulewidth}{0.3pt}
\renewcommand{\footrulewidth}{0pt}

% Couleurs
\definecolor{grisbox}{gray}{0.94}
\definecolor{bleutexte}{RGB}{30,60,120}

% Environnement de note
\newmdenv[backgroundcolor=grisbox, linewidth=0pt, innerleftmargin=10pt,
          innerrightmargin=10pt, innertopmargin=6pt, innerbottommargin=6pt,
          skipabove=6pt, skipbelow=6pt]{notebox}

% Style des sections
\usepackage{titlesec}
\titleformat{\section}[block]{\large\bfseries\color{bleutexte}}{}{0pt}{}[\vspace{-4pt}\rule{\linewidth}{0.4pt}]
\titlespacing{\section}{0pt}{14pt}{6pt}

% Numérotation des questions
\newcounter{qnum}
\newcommand{\question}[1]{%
  \stepcounter{qnum}%
  \medskip\noindent\textbf{Question \theqnum.}\quad #1\par
}

% Choix de réponse
\newcommand{\choix}[4]{%
  \begin{multicols}{2}
  \begin{itemize}[label={}, leftmargin=1.2em, itemsep=1pt, topsep=2pt]
    \item[\textbf{A)}] #1
    \item[\textbf{B)}] #2
    \item[\textbf{C)}] #3
    \item[\textbf{D)}] #4
  \end{itemize}
  \end{multicols}
}

% Réponse libre générique
\newcommand{\reponselibre}{%
  \par\smallskip
  \begin{notebox}
  \textit{Réponse :} \hfill\rule{8cm}{0.3pt}
  \vspace{4pt}

  \hspace*{3cm}\rule{9.4cm}{0.3pt}
  \end{notebox}
}

% Réponse avec attendus nommés : \reponseattendue{label1}{label2}
\newcommand{\reponseattendueA}[1]{%
  \par\smallskip
  \begin{notebox}
  \textit{Calcul :}\vspace{10pt}

  \hspace*{2cm}\rule{10.2cm}{0.3pt}\vspace{6pt}

  \hspace*{2cm}\rule{10.2cm}{0.3pt}\vspace{8pt}

  \noindent #1
  \end{notebox}
}

\begin{document}

% ─── TITRE ───────────────────────────────────────────────────────────────────
\begin{center}
  {\Large\bfseries\color{bleutexte} Contrôle de mathématiques}\\[4pt]
  {\large Équations différentielles \& Intégrales multiples}\\[6pt]
  \begin{tabular}{@{}r@{\hspace{4pt}}l@{\hspace{2cm}}r@{\hspace{4pt}}l@{}}
    \textbf{Durée :}   & 1 heure &
    \textbf{Barème :}  & QCM = 1 pt \quad Calcul = 2 pts \quad Total = 35 pts\\[2pt]
    \multicolumn{4}{l}{\textit{Calculatrice autorisée - Aucun autre document}}\\
  \end{tabular}
\end{center}

\vspace{2pt}
\noindent\rule{\linewidth}{0.6pt}
\vspace{4pt}

\begin{notebox}
\textbf{Nom :} \rule{6cm}{0.3pt}\quad
\textbf{Prénom :} \rule{6cm}{0.3pt}\\[6pt]
\textbf{Groupe :} \rule{4cm}{0.3pt}
\end{notebox}

% ─── PARTIE 1 ────────────────────────────────────────────────────────────────
\section{Équations différentielles du 1\textsuperscript{er} ordre}

\question{Laquelle de ces équations est une EDO linéaire du 1\textsuperscript{er} ordre ?}
\choix{$y'' + 2y = 0$}{$y' + 3y = e^x$}{$y \cdot y' = x$}{$(y')^2 = y$}

\question{La solution générale de $y' - 4y = 0$ est :}
\choix{$y = Ce^{4x}$}{$y = Ce^{-4x}$}{$y = 4x + C$}{$y = C\cos(4x)$}

\question{\textit{[Calcul]} On résout $y' + 2y = 0$ avec $y(0) = 5$.
Donner $y(x)$ puis calculer $y(1)$ (valeur approchée à $10^{-2}$).}
\reponseattendueA{$y(x) = $ \rule{4cm}{0.3pt} \hspace{1.5cm} $y(1) \approx $ \rule{3cm}{0.3pt}}

\question{La méthode de variation de la constante pour $y' + p(x)y = q(x)$ consiste à :}
\choix{Chercher $y_p = A \cdot q(x)$}
      {Poser $y = C(x) \cdot y_h$ où $y_h$ est la solution homogène}
      {Dériver l'équation pour éliminer $p(x)$}
      {Intégrer directement les deux membres}

\question{La solution générale de $y' + y = e^{2x}$ est de la forme :}
\choix{$y = Ce^{-x}$}
      {$y = Ce^{-x} + \tfrac{1}{3}e^{2x}$}
      {$y = Ce^{x} + e^{2x}$}
      {$y = Ce^{-x} + e^{2x}$}

% ─── PARTIE 2 ────────────────────────────────────────────────────────────────
\section{Équations différentielles du 2\textsuperscript{e} ordre}

\question{L'équation caractéristique de $y'' - 5y' + 6y = 0$ est :}
\choix{$r^2 - 5r + 6 = 0$}{$r^2 + 5r + 6 = 0$}{$r - 5 + 6 = 0$}{$r^2 - 5r = 0$}

\question{\textit{[Calcul]} Résoudre $r^2 - 5r + 6 = 0$. Donner $r_1$ et $r_2$.}
\reponseattendueA{$r_1 = $ \rule{3cm}{0.3pt} \hspace{2cm} $r_2 = $ \rule{3cm}{0.3pt}}

\question{D'après vos résultats, la solution générale de $y'' - 5y' + 6y = 0$ est :}
\choix{$y = C_1 e^{5x} + C_2 e^{6x}$}
      {$y = (C_1 + C_2 x)e^{3x}$}
      {$y = C_1 e^{2x} + C_2 e^{3x}$}
      {$y = e^{x}(C_1\cos(2x) + C_2\sin(3x))$}

\question{L'équation $y'' + 2y' + 5y = 0$ a pour discriminant $\Delta = -16$. Les racines sont :}
\choix{Réelles et distinctes}
      {Réelles égales}
      {Complexes conjuguées $-1 \pm 2i$}
      {Complexes conjuguées $1 \pm 2i$}

\question{La solution générale de $y'' + 2y' + 5y = 0$ est :}
\choix{$y = e^{-x}(C_1\cos 2x + C_2\sin 2x)$}
      {$y = e^{x}(C_1\cos 2x + C_2\sin 2x)$}
      {$y = (C_1 + C_2 x)e^{-x}$}
      {$y = C_1 e^{-x} + C_2 e^{5x}$}

\question{Pour $y'' - 4y = e^{2x}$, la forme $y_p = Ae^{2x}$ ne fonctionne pas car :}
\choix{$e^{2x}$ croît trop vite}
      {$e^{2x}$ est déjà solution de l'équation homogène}
      {Le second membre n'est pas polynomial}
      {L'équation est du 2\textsuperscript{e} ordre}

\question{On applique $y(0)=1$, $y'(0)=0$ à $y = C_1 e^{2x} + C_2 e^{3x}$. Le système obtenu est :}
\choix{$C_1+C_2 = 1$ \ et \ $2C_1+3C_2 = 0$}
      {$C_1+C_2 = 0$ \ et \ $2C_1+3C_2 = 1$}
      {$C_1 \cdot C_2 = 1$ \ et \ $C_1+C_2 = 0$}
      {$2C_1+3C_2=1$ \ et \ $C_1=C_2$}

\question{\textit{[Calcul]} Résoudre le système de la question précédente. Donner $C_1$ et $C_2$.}
\reponseattendueA{$C_1 = $ \rule{3cm}{0.3pt} \hspace{2cm} $C_2 = $ \rule{3cm}{0.3pt}}

% ─── PARTIE 3 ────────────────────────────────────────────────────────────────
\section{Intégrales doubles}

\question{\textit{[Calcul]} Calculer $\displaystyle\int_0^2\!\int_0^3 (x+y)\,dy\,dx$.}
\reponseattendueA{$\displaystyle\int_0^2\!\int_0^3 (x+y)\,dy\,dx \ = $ \rule{4cm}{0.3pt}}

\question{On calcule $\displaystyle\iint_D xy\,dA$ sur $D=[0,1]\times[0,2]$. Par séparation de Fubini :}
\choix{$\displaystyle\int_0^1 x\,dx \cdot \int_0^2 y\,dy$}
      {$\displaystyle\int_0^1\!\int_0^2 (x+y)\,dy\,dx$}
      {$\displaystyle\Bigl(\int_0^1 xy\,dx\Bigr)^2$}
      {$\displaystyle\int_0^2\!\int_0^1 xy\,dx\,dy$\quad\textit{(non séparé)}}

\question{\textit{[Calcul]} Calculer l'intégrale de la question précédente.}
\reponseattendueA{$\displaystyle\iint_D xy\,dA \ = $ \rule{4cm}{0.3pt}}

\question{Pour $\displaystyle\iint_D \sqrt{x^2+y^2}\,dA$ sur $x^2+y^2\leq 9$, le meilleur changement de variables est :}
\choix{$x=u+v,\ y=u-v$}
      {Coordonnées polaires : $x=r\cos\theta,\ y=r\sin\theta$}
      {$x=r\cosh\theta,\ y=r\sinh\theta$}
      {Aucun changement nécessaire}

\question{\textit{[Calcul]} Avec les coordonnées polaires, calculer $\displaystyle\iint_D \sqrt{x^2+y^2}\,dA$.\\
\textit{Rappel :} $\sqrt{x^2+y^2}=r$, $dA = r\,dr\,d\theta$.}
\reponseattendueA{$\displaystyle\iint_D \sqrt{x^2+y^2}\,dA \ = $ \rule{4cm}{0.3pt}}

\question{On inverse l'ordre dans $\displaystyle\int_0^1\!\int_x^1 f(x,y)\,dy\,dx$. L'intégrale équivalente est :}
\choix{$\displaystyle\int_0^1\!\int_0^y f(x,y)\,dx\,dy$}
      {$\displaystyle\int_0^1\!\int_0^1 f(x,y)\,dx\,dy$}
      {$\displaystyle\int_0^1\!\int_y^1 f(x,y)\,dx\,dy$}
      {$\displaystyle\int_0^1\!\int_0^x f(x,y)\,dx\,dy$}

% ─── PARTIE 4 ────────────────────────────────────────────────────────────────
\section{Green-Riemann \& Intégrales triples}

\question{La formule de Green-Riemann s'applique si :}
\choix{$D$ est un domaine fermé borné à bord régulier orienté positivement}
      {$f$ est continue sur $\mathbb{R}^2$}
      {$P$ et $Q$ sont des constantes}
      {Le domaine est un rectangle uniquement}

\question{\textit{[Calcul]} Avec $P=-y$, $Q=x$, calculer l'aire du disque de rayon $R=3$ via Green-Riemann.\\
\textit{Rappel :} $\oint_{\partial D}(-y\,dx + x\,dy) = 2\,\mathrm{Aire}(D)$.}
\reponseattendueA{$\mathrm{Aire}(D) = $ \rule{4cm}{0.3pt}}

\question{Le jacobien en coordonnées sphériques $(\rho,\varphi,\theta)$ vaut :}
\choix{$\rho^2$}{$\rho\sin\varphi$}{$\rho^2\sin\varphi$}{$\rho^2\cos\varphi$}

\question{\textit{[Calcul]} Calculer le volume de la boule $x^2+y^2+z^2\leq 1$ en coordonnées sphériques.\\
\textit{Rappel :} $\displaystyle\int_0^\pi\sin\varphi\,d\varphi = 2$.}
\reponseattendueA{$V = \displaystyle\iiint_V dV = $ \rule{4cm}{0.3pt}}

% ─── PARTIE 5 ────────────────────────────────────────────────────────────────
\section{Applications}

\question{La coordonnée $\bar{x}$ du centre de gravité d'une surface $D$ d'aire $A$ (densité $\rho=1$) est :}
\choix{$\bar{x} = \iint_D x^2\,dA$}
      {$\bar{x} = A \cdot \iint_D x\,dA$}
      {$\bar{x} = \dfrac{1}{A}\iint_D x\,dA$}
      {$\bar{x} = \iint_D x\,dA$}

\question{\textit{[Calcul]} Calculer $\bar{x}$ du triangle de sommets $(0,0)$, $(3,0)$, $(0,3)$ (aire $= 9/2$).\\
\textit{Rappel :} $\displaystyle\iint_D x\,dA = \int_0^3\!\int_0^{3-x} x\,dy\,dx$.}
\reponseattendueA{$\displaystyle\iint_D x\,dA = $ \rule{3.5cm}{0.3pt} \hspace{1.5cm} $\bar{x} = $ \rule{3cm}{0.3pt}}

% ─── CORRIGÉ ─────────────────────────────────────────────────────────────────
\newpage
\begin{center}
  {\large\bfseries\color{bleutexte} Corrigé détaillé}\\[2pt]
  {\small\itshape Les réponses QCM sont encadrées. Les calculs sont détaillés étape par étape.}
\end{center}

\vspace{2pt}

% ── Grille QCM ───────────────────────────────────────────────────────────────
\begin{notebox}
\textbf{Réponses QCM :}\quad
Q1\,\textbf{B} \quad Q2\,\textbf{A} \quad Q4\,\textbf{B} \quad Q5\,\textbf{B} \quad
Q6\,\textbf{A} \quad Q8\,\textbf{C} \quad Q9\,\textbf{C} \quad Q10\,\textbf{A} \quad
Q11\,\textbf{B} \quad Q12\,\textbf{A} \quad Q15\,\textbf{A} \quad Q17\,\textbf{B} \quad
Q19\,\textbf{A} \quad Q20\,\textbf{A} \quad Q22\,\textbf{C} \quad Q24\,\textbf{C}
\end{notebox}

\vspace{4pt}

% ── Justifications QCM ───────────────────────────────────────────────────────
\noindent\textbf{Justifications QCM}

\smallskip
\noindent\textbf{Q1 — B.} Une EDO \emph{linéaire} du 1\ier{} ordre s'écrit $y' + p(x)y = q(x)$ : les termes $y$ et $y'$ apparaissent à la puissance 1, sans produit entre eux. Les propositions C et D sont non linéaires ($y \cdot y'$ ou $(y')^2$) ; A est du 2\ieme{} ordre.

\smallskip
\noindent\textbf{Q2 — A.} L'équation caractéristique de $y' - 4y = 0$ donne $r = 4$, donc $y = Ce^{4x}$.

\smallskip
\noindent\textbf{Q4 — B.} La variation de la constante remplace la constante $C$ par une \emph{fonction} $C(x)$ dans la solution homogène : $y = C(x)\cdot y_h$. On dérive et on substitue pour trouver $C'(x)$.

\smallskip
\noindent\textbf{Q5 — B.} Pour $y'+y=e^{2x}$ : $y_h=Ce^{-x}$. On cherche $y_p = Ae^{2x}$ $\Rightarrow$ $y_p'=2Ae^{2x}$ $\Rightarrow$ $2Ae^{2x}+Ae^{2x}=e^{2x}$ $\Rightarrow$ $3A=1$ $\Rightarrow$ $A=\tfrac{1}{3}$.

\smallskip
\noindent\textbf{Q6 — A.} On substitue $y''\to r^2$, $y'\to r$, $y\to 1$ : $r^2-5r+6=0$.

\smallskip
\noindent\textbf{Q8, Q9, Q10.} Voir calculs détaillés Q7 et Q13 ci-dessous.

\smallskip
\noindent\textbf{Q11 — B.} $e^{2x}$ est solution de l'homogène $y''-4y=0$ (car $r=2$ est racine). La forme $y_p=Ae^{2x}$ serait absorbée par $y_h$. On prend $y_p=Axe^{2x}$.

\smallskip
\noindent\textbf{Q15 — A.} $f(x,y)=xy$ est \emph{séparable} ($g(x)=x$, $h(y)=y$), sur un rectangle $\Rightarrow$ Fubini : $\int_0^1 x\,dx \cdot \int_0^2 y\,dy$.

\smallskip
\noindent\textbf{Q17 — B.} La symétrie circulaire du disque et la présence de $\sqrt{x^2+y^2}$ imposent les coordonnées polaires.

\smallskip
\noindent\textbf{Q19 — A.} Domaine $\{(x,y):0\le x\le 1,\ x\le y\le 1\}$ = triangle. En inversant : $y$ varie de 0 à 1, et pour chaque $y$, $x$ varie de 0 à $y$.

\smallskip
\noindent\textbf{Q20 — A.} Green-Riemann nécessite un domaine $D$ fermé et borné, à bord $\partial D$ régulier et orienté \emph{positivement} (sens trigonométrique).

\smallskip
\noindent\textbf{Q22 — C.} En sphériques : $x=\rho\sin\varphi\cos\theta$, $y=\rho\sin\varphi\sin\theta$, $z=\rho\cos\varphi$. Le jacobien de la transformation est $\rho^2\sin\varphi$.

\smallskip
\noindent\textbf{Q24 — C.} Par définition : $\bar{x} = \dfrac{1}{A}\iint_D x\,dA$ (barycentre pondéré par l'aire, densité uniforme).

\vspace{6pt}

% ── Calculs détaillés ────────────────────────────────────────────────────────
\noindent\textbf{Calculs détaillés}

\smallskip
\noindent\textbf{Q3.} $y'+2y=0$ est homogène. L'équation caractéristique $r+2=0$ donne $r=-2$.
$$y(x) = Ce^{-2x}$$
Condition initiale $y(0)=5$ : $C\cdot e^0 = 5 \Rightarrow C=5$.
$$\boxed{y(x) = 5e^{-2x}} \qquad y(1) = 5e^{-2} \approx 5\times 0{,}1353 \approx \mathbf{0{,}68}$$

\smallskip
\noindent\textbf{Q7.} $r^2-5r+6=0$. Discriminant $\Delta = 25-4\times6 = 1 > 0$.
$$r = \frac{5\pm 1}{2} \qquad \Rightarrow \qquad \boxed{r_1 = 2,\quad r_2 = 3}$$
Deux racines réelles distinctes $\Rightarrow$ solution : $y = C_1e^{2x}+C_2e^{3x}$ \textbf{(Q8 : C)}.

\smallskip
\noindent\textbf{Q13.} Système issu de $y(0)=1$, $y'(0)=0$ avec $y=C_1e^{2x}+C_2e^{3x}$ :
$$\begin{cases}C_1+C_2 = 1\\ 2C_1+3C_2 = 0\end{cases}$$
De la 2\ieme{} : $C_2 = -\tfrac{2}{3}C_1$. On substitue : $C_1 - \tfrac{2}{3}C_1 = 1 \Rightarrow \tfrac{1}{3}C_1=1$.
$$\boxed{C_1 = 3,\quad C_2 = -2}$$

\smallskip
\noindent\textbf{Q14.} On intègre d'abord en $y$ (Fubini) :
$$\int_0^3(x+y)\,dy = \Bigl[xy+\tfrac{y^2}{2}\Bigr]_0^3 = 3x+\tfrac{9}{2}$$
Puis en $x$ :
$$\int_0^2\!\Bigl(3x+\tfrac{9}{2}\Bigr)dx = \Bigl[\tfrac{3x^2}{2}+\tfrac{9x}{2}\Bigr]_0^2 = 6+9 = \boxed{18}$$

\smallskip
\noindent\textbf{Q16.} Séparation de Fubini ($f=xy$ séparable) :
$$\int_0^1 x\,dx \cdot \int_0^2 y\,dy = \Bigl[\tfrac{x^2}{2}\Bigr]_0^1 \times \Bigl[\tfrac{y^2}{2}\Bigr]_0^2 = \tfrac{1}{2}\times 2 = \boxed{1}$$

\smallskip
\noindent\textbf{Q18.} En polaires sur $x^2+y^2\le 9$ : $r\in[0,3]$, $\theta\in[0,2\pi]$, $\sqrt{x^2+y^2}=r$, $dA=r\,dr\,d\theta$.
$$\iint_D\sqrt{x^2+y^2}\,dA = \int_0^{2\pi}\!\!d\theta\int_0^3 r\cdot r\,dr = 2\pi\int_0^3 r^2\,dr = 2\pi\Bigl[\tfrac{r^3}{3}\Bigr]_0^3 = 2\pi\cdot 9 = \boxed{18\pi}$$

\smallskip
\noindent\textbf{Q21.} On paramètre le cercle $\partial D$ de rayon 3 : $x=3\cos t$, $y=3\sin t$, $t\in[0,2\pi]$.
$$\oint(-y\,dx+x\,dy) = \int_0^{2\pi}(9\sin^2 t + 9\cos^2 t)\,dt = 9\int_0^{2\pi}dt = 18\pi$$
Donc $\mathrm{Aire}(D) = \tfrac{18\pi}{2} = \boxed{9\pi}$ \quad (cohérent avec $\pi R^2 = 9\pi$).

\smallskip
\noindent\textbf{Q23.} En sphériques, $dV=\rho^2\sin\varphi\,d\rho\,d\varphi\,d\theta$, bornes : $\rho\in[0,1]$, $\varphi\in[0,\pi]$, $\theta\in[0,2\pi]$.
$$V = \int_0^{2\pi}\!\!d\theta\int_0^\pi\!\sin\varphi\,d\varphi\int_0^1\rho^2\,d\rho = 2\pi \times 2 \times \tfrac{1}{3} = \boxed{\dfrac{4\pi}{3}}$$

\smallskip
\noindent\textbf{Q25.} On calcule d'abord $\iint_D x\,dA$ sur le triangle ($y$ de $0$ à $3-x$) :
$$\int_0^3\!\int_0^{3-x}x\,dy\,dx = \int_0^3 x(3-x)\,dx = \Bigl[{\tfrac{3x^2}{2}-\tfrac{x^3}{3}}\Bigr]_0^3 = \tfrac{27}{2}-9 = \tfrac{9}{2}$$
Puis : $\bar{x} = \dfrac{\iint_D x\,dA}{A} = \dfrac{9/2}{9/2} = \boxed{1}$

\end{document}
